% Generated draft: requires mathematical and source review.
\documentclass[11pt]{article}
\usepackage[margin=25mm]{geometry}
\usepackage{fontspec}
\setmainfont{Latin Modern Roman}
\usepackage{amsmath,amssymb,mathtools,unicode-math}
\setmathfont{Latin Modern Math}
\usepackage{enumitem,needspace,xurl,hyperref}
\setlength{\parindent}{0pt}
\setlength{\parskip}{5pt}
\setlength{\emergencystretch}{3em}
\newcommand{\N}{\mathbb{N}}
\newcommand{\Z}{\mathbb{Z}}
\newcommand{\Q}{\mathbb{Q}}
\newcommand{\R}{\mathbb{R}}
\newcommand{\C}{\mathbb{C}}
\newcommand{\F}{\mathbb{F}}
\newcommand{\A}{\mathbb{A}}
\newcommand{\PP}{\mathbb{P}}
\newcommand{\E}{\mathbb{E}}
\newcommand{\eps}{\varepsilon}
\newcommand{\dd}{\,\mathrm{d}}
\begin{document}
\section*{UnsolvedMath ID: 801137}
\textbf{Source problem code:} catalogue entry
\subsection{Definitions and mathematical statement}
In the type \(\mathrm{II}_1\) setting indicated by the source title, let \(N\) be a factor with normalized faithful normal trace \(\tau\), and let \(\mathcal U\) be an ultrafilter on an index set \(I\). Define \(\|x\|_{2,\tau}=\tau(x^*x)^{1/2}\) and the tracial ultrapower
\[
N^{\mathcal U}=\ell^\infty(I,N)\big/\{(x_i)_{i\in I}:\lim_{i\to\mathcal U}\|x_i\|_{2,\tau}=0\}.
\]
Here \(\ell^\infty(I,N)\) consists of the operator-norm-bounded families. A subfactor is a unital von Neumann subalgebra that is itself a factor, and \([M:P]\) denotes the Jones index of a subfactor inclusion \(P\subset M\).

With the ambient parameters fixed, the source asks whether every subfactor \(N_0\subset N^{\mathcal U}\) satisfying
\[
[N^{\mathcal U}:N_0]=2
\]
admits a family of subfactors \((N_i\subset N)_{i\in I}\) such that
\[
[N:N_i]=2\quad\text{for every }i\in I,
\qquad
N_0\cong\prod_{i\to\mathcal U}(N_i,\tau|_{N_i}).
\]
The product on the right is the tracial ultraproduct, defined by the analogous quotient of bounded families using the restricted traces. The symbol \(\cong\) denotes an isomorphism of von Neumann algebras; the supplied statement does not require this isomorphism to respect the inclusions into \(N^{\mathcal U}\).

The supplied statement does not explicitly specify the hypotheses on \(N\), the index set \(I\), or whether \(\mathcal U\) must be nonprincipal. The type \(\mathrm{II}_1\) interpretation comes from the source title; any further intended restrictions, including separability or countability, remain unresolved.
\subsection{Short English statement}
Must every Jones-index-\(2\) subfactor of a tracial ultrapower of a type \(\mathrm{II}_1\) factor be isomorphic to a tracial ultraproduct of index-\(2\) subfactors of the original factor? The supplied record leaves the precise ambient and ultrafilter hypotheses unspecified.
\subsection{Sources}
\begin{itemize}
\item UnsolvedMath local registry, permanent ID \texttt{801137}, source problem code \texttt{catalogue entry}, upstream revision \texttt{d371ee917ac9647fbb06cd6fc529deec6f3dda82}. The registry attributes the problem to Elayavalli (2025) and labels it open; these are supplied metadata, not an independently verified status assessment.
\item Elayavalli, \emph{50 Open Problems: Ultraproduct \(\mathrm{II}_1\) Factors} (2025), Problem 38. Supplied source URL: \url{https://arxiv.org/abs/2511.20377}. The record reports access on 2026-10-16 and extraction by \texttt{arxiv\_latex}; these are supplied provenance statements, not a claim of independent source consultation.
\end{itemize}
\end{document}
